\documentclass[../CLT_SYK.tex]{subfiles}
\setcounter{section}{0}

\begin{document}

\section{Introduction and main results}
\label{sec:introduction}



\subsection{The SYK model}\label{subsec:SYK-model}
The Sachdev--Ye--Kitaev (SYK) model is a quantum-mechanical model of \(N\) Majorana fermions with random many-body interactions. It originates in the work of Sachdev and Ye on random quantum magnets \cite{SachdevYe1993} and in Kitaev’s subsequent formulation in terms of Majorana fermions \cite{Kitaev2015}. Later work by Maldacena, Stanford, and others clarified its large-\(N\) behavior and revealed deep connections to quantum chaos and two-dimensional quantum gravity \cite{MaldacenaStanford2016}. These features make the SYK model a useful setting for studying strongly interacting quantum many-body systems, quantum chaos, and quantum gravity.

Depending on how the inverse temperature \(\beta=\beta_N\) scales with \(N\), the model exhibits several distinct regimes. For fixed \(\beta\), one obtains the standard finite-temperature large-\(N\) regime; when \(1\ll \beta\ll N\), the system enters the conformal regime, characterized by strong interactions and non-Fermi liquid behavior. At lower temperatures, \(N\ll \beta\), the quantum-Schwarzian regime emerges, in which finite-\(N\) effects and quantum fluctuations of the Schwarzian mode play a central role. Finally, when \(\beta\) is exponentially large in \(N\), the partition function is mainly determined by the eigenvalues near the edge of the spectrum. In this regime, random matrix theory  describes their fluctuations. The appearance of such diverse crossover phenomena is rare in a strongly interacting model that remains analytically tractable in suitable limits.

Previous mathematical work has studied the limiting free energy of the SYK model for fixed $\beta$ \cite{GPSZ2026,He2026}. Moreover, Anschuetz, Gamarnik, and Kiani \cite{AGK2025} established a universality result showing that the limiting free energy is the same for a broad class of interaction distributions.  This naturally raises the question of how physical quantities depend on the interactions. Motivated by this question, we study fluctuations of the free energy and show that, for sufficiently small \(|\beta|\), the suitably centered and rescaled free energy converges to a Gaussian distribution.

Fix an even integer $q\geq4$, and let $N$ be even. Set $[N] \coloneqq \{1,2,\ldots,N\}$ and write
\[
  \calI_N \coloneqq \{I\subset[N]:|I|=q\},\qquad
  \sigma_N^2 \coloneqq \frac{(q-1)!}{N^{q-1}},
\]
where $|I|$ denotes the cardinality of the set $I$.
Let $G \coloneqq (g_I)_{I\in\calI_N}$ be a family of independent standard Gaussian random
variables. We denote the corresponding probability measure and expectation by $\mathbb P$ and $\E$, respectively. Following \cite[Section~2]{GPSZ2026}, we let $\psi_1,\ldots,\psi_N$ be Majorana matrices of size $2^{N/2} \times 2^{N/2}$, which satisfy the following anticommutation relation\footnote{\textcolor{red}{Throughout the paper, use a consistent word ``anticommutation relation''}}:
\begin{equation}
  \label{eq: clifford relation}
  \{\psi_i,\psi_j\} \coloneqq \psi_i \psi_j + \psi_j \psi_i =2\delta_{ij}\I,
\end{equation}
where $\I$ denotes the identity matrix.
(NB. The Majorana matrices considered here are $\sqrt{2}$ times those used in \cite{GPSZ2026,He2026}; this is merely a difference in normalization.)
For
$I = \{i_1, \ldots, i_q\} \in \calI_N$ with $i_1 < \cdots < i_q$, set
\[
  \psi_I \coloneqq \mathrm i^{q/2}\psi_{i_1}\cdots\psi_{i_q},
\]
where $\mathrm i$ denotes the imaginary unit. Since Majorana matrices are Hermitian, it follows from \eqref{eq: clifford relation} that 
$\psi_I$ is also Hermitian. Moreover, $\psi_I^2=\I$ and, in particular, $\norm{\psi_I}_\infty=1$, where $\norm{\,\cdot\,}_\infty$ denotes the operator norm of a matrix.  
The Hamiltonian of the SYK model is defined by 
\[
  H_N(G) \coloneqq \sigma_N\sum_{I\in\calI_N}g_I\psi_I.
\]
The corresponding partition function is then given by
\[
  Z_N(\beta) \coloneqq Z_N(\beta;G) \coloneqq \tr e^{-\beta H_N(G)},\qquad \beta \in \mathbb R.
\]
Here, $\tr \coloneqq 2^{-N/2}\operatorname{Tr}$ denotes the normalized trace; in
particular, $Z_N(0;G)=1$.  

\subsection{Main results}
Here, we state our main results: Theorems~\ref{thm:partition-clt} and~\ref{thm:free-energy-clt}.
We first recall the convergence of the annealed free energy, which was established for sufficiently small $|\beta|$ in \cite{GPSZ2026} and recently extended to all $\beta$ in \cite{He2026}. Moreover, in \cite{He2026}, the limit is identified with the Schwinger--Dyson pressure. Throughout this article, all limits as $N \to \infty$ are taken through even integers.
\begin{theorem}[{\cite[Theorem~1.1 and Proposition~3.3]{He2026}}]\label{thm: input}
For any $\beta\in \mathbb R$, there exists  $\ell_q(\beta)$ such that
\begin{align*}
    \lim_{N\to \infty} N^{-1} \log \E[Z_N] =\ell_q(\beta).
\end{align*}
Moreover, $\ell_q(\beta)$ is differentiable with respect to $\beta$ on $\mathbb R$.
\end{theorem}

\textcolor{red}{We also prove that $\beta \ell_q'(\beta)>0$ for every $\beta\in\mathbb R\setminus\{0\}$ (see Proposition~\ref{prop:annealed-free-energy}).}


Our first main result concerns the fluctuations of the partition function ${Z_N(\beta)}$. For $a > 0$, we write $\N(0, a)$ for the centered Gaussian distribution with variance $a$.
\begin{theorem}
\label{thm:partition-clt}
There exists $\beta_q\in (0,\infty)$ such that for any $\beta\in (-\beta_q,\beta_q) \setminus \{0\}$,
 \begin{equation}\label{eq:scaled-relative-clt}
  N^{q/2-1}\left( \frac{Z_N(\beta)}{\E[Z_N(\beta)]} - 1\right)
  \Longrightarrow
  \,\N\left(0,\frac{q!}{2}\beta^2(\ell_q'(\beta))^2\right),
\end{equation}
where the arrow $\Longrightarrow$ means convergence in distribution. Moreover,
\begin{equation}\label{eq:standardized-clt}
  \frac{Z_N(\beta)-\E[Z_N(\beta)]}
  {\sqrt{\Var Z_N(\beta)}}
  \Longrightarrow\N(0,1),
\end{equation}
where $\Var Z_N(\beta)$ denotes the variance of $Z_N(\beta)$.
\end{theorem}

Combining the above theorem with the asymptotic relation $\log x \sim x-1$ as $x \to 1$, we obtain the following result on the fluctuations of the free energy.
\begin{theorem}\label{thm:free-energy-clt}
With the same $\beta_q$ as in Theorem~\ref{thm:partition-clt}, for any $\beta\in (-\beta_q,\beta_q) \setminus \{0\}$,
\[
  N^{q/2-1}\left(
    \log Z_N(\beta)-\E[\log Z_N(\beta)]
  \right)
  \Longrightarrow
  \,\N\left(0,\frac{q!}{2}\,\beta^2(\ell_q'(\beta))^2\right).
\]
\end{theorem}

\begin{remark}
The proof of Theorem~\ref{thm:partition-clt} is based on a polynomial chaos expansion of $Z_N(\beta)$, and the smallness of $|\beta|$ is used to control the higher-order chaos terms; see Proposition~\ref{prop:higher-chaos-closure} below.
We do not pursue an extension to all $\beta\in\mathbb R$ here, nor do we make a conjecture as to whether the same fluctuation results hold at all temperatures.
\end{remark}

\begin{remark}
In related work, Feng, Tian, and Wei \cite{FTW2018} proved a central limit theorem for linear eigenvalue statistics of the Hamiltonian $H_N$. Their result applies to a fixed test function \(f\), under suitable regularity assumptions on \(f\) and \(f'\). In their normalization, however, the partition function corresponds to
\[
f_{N,\beta}(x) \coloneqq e^{-\beta \sqrt{N} x},
\]
so both the test function and its derivative depend on \(N\). Therefore, their result does not apply directly to \(Z_N(\beta)\).
\end{remark}


\subsection{Outline of the proofs and organization of the paper}

We give an outline of the proofs of the main results.
Define the normalized partition function by
\begin{equation} \label{eq:Y_N}
  Y_N(\beta) \coloneqq Y_N(\beta; G) \coloneqq {Z_N(\beta)}/{\E[Z_N(\beta)]}.
\end{equation}
By applying a polynomial chaos expansion to $Y_N(\beta)$, we write 
\[
  Y_N(\beta) = \sum_{k = 0}^\infty P_kY_N(\beta),
\]
where $P_k Y_N(\beta)$ denotes the $k$-th chaos component.
(The precise definition of the above expansion is given in Section~\ref{sec:chaos-decomposition}.) The normalization gives $P_0 Y_N(\beta) = \E[Y_N(\beta)] = 1$. Moreover, in Lemma~\ref{lem:exact-second-chaos}, using the anticommutation relation~\eqref{eq: clifford relation} and Gaussian integration by parts, we deduce that $P_1 Y_N(\beta) = 0$ and
\[ 
  P_2Y_N(\beta) = A_N(\beta) Q_N,
\]
where
\begin{equation}\label{eq:finite-N-FTW-coefficient}
      A_N(\beta) \coloneqq \frac{\beta}{2}\partial_\beta\log \E[Z_N(\beta)],
  \qquad
    Q_N \coloneqq \frac1{|\mathcal I_N|}\sum_{I\in\mathcal I_N}(g_I^2-1).
\end{equation}
Hence, we obtain
\begin{equation}\label{eq:intro-chaos-expansion}
  Y_N(\beta)-1=
  A_N(\beta)Q_N+\mathcal R_N(\beta),
\end{equation}
where $\mathcal R_N(\beta) \coloneqq \sum_{k \geq 3} P_k Y_N(\beta)$.
Theorem~\ref{thm: input} and the convexity of $\log \E[Z_N(\beta)]$ in $\beta$ imply that
\begin{equation}\label{eq:intro-A_N-limit}
  \frac{A_N(\beta)}N
  = 
  \frac{\beta}{2}\cdot \frac{1}{N}\partial_\beta\log \E[Z_N(\beta)]
  \xlongrightarrow{N \to \infty} \frac{\beta}{2}\ell_q'(\beta),
\end{equation}
as shown in Proposition~\ref{prop:annealed-free-energy}.
Since the usual central limit theorem applies to $Q_N$, Theorem~\ref{thm:partition-clt} follows from \eqref{eq:intro-chaos-expansion} and \eqref{eq:intro-A_N-limit}, combined with  
\begin{equation}\label{eq:intro-higher-chaoses}
  \E \bigl[ \mathcal R_N(\beta)^2 \bigr]= o(N^{-(q-2)}),\qquad \text{as } N \to \infty,
\end{equation}
where $o(\cdot)$ denotes the usual Landau notation.
The above estimate is proven in Proposition~\ref{prop:higher-chaos-closure} 
by using a combinatorial argument to bound the number of non-degenerate terms in the higher-order chaoses\footnote{\textcolor{red}{Please check the expression ``non-degenerate terms''. Do you agree with this expression?}}.

To obtain Theorem~\ref{thm:free-energy-clt}, observe that
\begin{align*}
  &N^{q/2-1}\bigl(\log Z_N(\beta)-\E[\log Z_N(\beta)]\bigr)\\
  &=N^{q/2-1}\bigl(\log Y_N(\beta)-(Y_N(\beta)-1)\bigr)
   +N^{q/2-1}(Y_N(\beta)-1)-N^{q/2-1}\E[\log Y_N(\beta)].
\end{align*}
The Taylor expansion $\log(1+x)-x=O(x^2)$ as $x \to 0$, together with Theorem~\ref{thm:partition-clt}, shows that the first term  on the right-hand side converges to zero in probability. Moreover, using the estimates obtained in the proof of \eqref{eq:intro-higher-chaoses}, we deduce that $\E[(Y_N(\beta)-1)^2]=O(N^{2-q})$ and
$\mathbb P(Y_N(\beta)<1/2)=O(N^{-3(q/2-1)})$. Together with
$Z_N(\beta)\geq1$, these estimates give
\[
  \E[\log Y_N(\beta)]=O\!\left(N^{2-q}\right)
  =o\!\left(N^{1-q/2}\right).
\]
This completes the proof of Theorem~\ref{thm:free-energy-clt}; see Section~\ref{sec:main-results} for details.

\textcolor{red}{The rest of the paper is organized as follows. In
Section~\ref{sec:chaos-decomposition}, we introduce the polynomial chaos expansion
of $Z_N(\beta)$ and compute its first two chaos components. We show that the first
chaos vanishes and that the second chaos equals $A_N(\beta)Q_N$, where $A_N(\beta)$ and $Q_N$ 
are defined in~\eqref{eq:finite-N-FTW-coefficient}. In
Section~\ref{sec:Rademacher-remainder}, we show that the chaos components of order
three and higher are negligible compared with the second chaos. In
Section~\ref{sec:main-results}, we derive the asymptotics of $A_N(\beta)$ from the
annealed free energy estimates of~\cite{He2026}. Combining these results, we prove
Theorems~\ref{thm:partition-clt} and~\ref{thm:free-energy-clt}.}
\end{document}
